\documentclass[11pt,a4paper]{article}
\usepackage[T1]{fontenc}
\usepackage{lmodern,amsmath,amssymb,amsthm,microtype,booktabs}
\usepackage[margin=25mm]{geometry}
\usepackage[colorlinks=true,urlcolor=blue,linkcolor=blue]{hyperref}
\newcommand{\dd}{\mathop{}\!\mathrm d}
\newcommand{\sech}{\operatorname{sech}}
\newcommand{\cD}{\mathcal D}
\newcommand{\cS}{\mathcal S}
\title{Local scalar insertions in dS$_3$:\\from normalized BD modes to global CG coefficients}
\author{Alan Day \and Chengheng\thanks{AI research assistant (OpenAI).}}
\date{Working research note --- 31 October 2026}
\begin{document}\maketitle
\begin{abstract}
For an explicitly chosen $G=0$ local multiplication insertion, we
compute the scalar trilinear coefficient using Klein--Gordon normalized
Bunch--Davies modes in global dS$_3$. A finite Beta sum evaluates
arbitrary finite angular-mode entries. Future boundary data fix the
unitary shadow phases, and a normalized spherical period fixes the
full complex global Clebsch--Gordan coefficient. Its norm differs from
the previously extracted paired $C_0$ factor by explicit leg factors.
This is a local dictionary in a declared model, not a construction of
the physical observer mass measure or finite-gravity sewing.
\end{abstract}

\section{The vacuum and its normalized modes}
Set the radius to one and take
\begin{equation}
 ds^2=-dt^2+\cosh^2t\,d\Omega_2^2,\qquad
 m_{\rm eff}^2=1+\nu^2,\quad \nu>0.
\end{equation}
Use $\int|Y_{lm}|^2\dd\Omega=1$ and $Y_{00}=1/\sqrt{4\pi}$.
The angular momentum $l$ is an SO(3) $K$-type label.
A single scalar principal-series representation, with global spin
$J=0$, contains \emph{every} $l\geq0$.

For $u_{\nu lm}=f_{\nu l}Y_{lm}$ and $\chi_{\nu l}=\cosh t\,f_{\nu l}$,
\begin{equation}
 \chi''+[\nu^2+l(l+1)\sech^2t]\chi=0,\qquad
 i(\chi^*\chi'-\chi'^*\chi)=1.
 \label{eq:ode}
\end{equation}
LMW \cite{lmw} distinguish the reflectionless in/out modes from the
Euclidean or BD modes. Regularity on the lower Euclidean hemisphere
and KG normalization give
\begin{equation}
 \chi_{\nu0}=\frac{\sinh(\pi\nu/2-i\nu t)}{\sqrt{\nu\sinh\pi\nu}}
 =N_\nu(e^{-i\nu t}-q_\nu e^{i\nu t}),\quad
 q_\nu=e^{-\pi\nu},\quad N_\nu=[2\nu(1-q_\nu^2)]^{-1/2}.
 \label{eq:bd}
\end{equation}
The numerator vanishes at $t=-i\pi/2$, making $f_{\nu0}$ regular.
Omitting the second exponential changes the vacuum and the vertex.

\newpage
\section{An actual local insertion and a convergent mode integral}
Define $D_j=\partial_t-j\tanh t$ and
$d_l(\nu)=\prod_{j=1}^l\sqrt{\nu^2+j^2}$. The normalized modes are
\begin{equation}
 f_{\nu l}=\frac{i^l}{d_l(\nu)\cosh t}
 D_l\cdots D_1\chi_{\nu0}.
 \label{eq:darboux}
\end{equation}
Each Darboux step maps the previous equation into \eqref{eq:ode},
multiplies its Wronskian by $\nu^2+j^2$ and preserves the regular
Euclidean branch. The phase $i^l$ agrees with MM (16)--(17)
\cite{mm}:
\begin{equation}
 f_{\nu l}=
 \sqrt{\frac{\Gamma(l+1-i\nu)\Gamma(l+1+i\nu)}2}\,
 (\cosh t)^{-1/2}P_{-1/2+i\nu}^{-l-1/2}(i\sinh t).
 \label{eq:legendre}
\end{equation}
The type-2 Legendre branch is intended.

Following the structure of KL (5.4) \cite{kl}, we \emph{define}
the present dS$_3$ model by
\begin{equation}
 \langle o|\mathcal O_\phi|i\rangle
 :=\int_{\mathrm{dS}_3}\dd V\,u_o^*(X)\phi(X)u_i(X).
 \label{eq:insertion}
\end{equation}
Here $\phi=\sum(a_pu_p+a_p^\dagger u_p^*)$ is a canonically
normalized free BD scalar. The local coupling is chosen to be one.
Thus the annihilation and emission coefficients are
$\int u_o^*u_i u_p\,\dd V$ and $\int u_o^*u_i u_p^*\,\dd V$.
The construction is currently at fixed representation labels; a
physical mass direct integral and operator domain are additional tasks.

Write $T_\eta=\int\dd V\,\prod_j u_j^{(\eta_j)}$, where $\eta_j=+1$
means $u_j$ and $\eta_j=-1$ means $u_j^*$. For finite angular modes,
$u_j=O(e^{-|t|})$, so the integrand with volume is $O(e^{-|t|})$.
The integral converges absolutely. Killing derivatives give a
total divergence with vanishing boundary flux, establishing
invariance. Finite group translates retain this decay because the
induced conformal factors on the compact boundary sphere are bounded.

For explicit evaluation, set $y=\tanh t$ and
\begin{equation}
 P_{0,\sigma}=1,\quad
 P_{l,\sigma}=(1-y^2)P'_{l-1,\sigma}
 +(i\sigma\nu-ly)P_{l-1,\sigma},\qquad \sigma=\pm1.
\end{equation}
There are eight exponential branches in a triple product.
The saved evaluator uses $m_j=0$; general magnetic labels enter through
the usual conjugated Gaunt coefficient. Every radial term uses
\begin{equation}
 I_n(\omega)=\int_{\mathbb R}\sech t\,\tanh^nt\,e^{i\omega t}\dd t
 =\sum_{k=0}^n{n\choose k}2^k(-1)^{n-k}
 B\left(k+\frac{1+i\omega}{2},\frac{1-i\omega}{2}\right).
 \label{eq:beta}
\end{equation}
This follows by $x=(1+\tanh t)/2$; $I_0=\pi/\cosh(\pi\omega/2)$.
If $b_j$ sums the absolute polynomial coefficients of both
normalized branches, truncation to $|t|\leq T$ has
\begin{equation}
 |E_T|\leq4e^{-T}|\mathcal G|\prod_j b_j.
 \label{eq:tail}
\end{equation}
This bounds the tail, not the numerical integration error inside it.

\newpage
\section{The spherical entry, future data and shadow phases}
For $S=\nu_1+\nu_2+\nu_3$, introduce
\begin{equation}
 \cD=\cosh\frac{\pi S}{2}\prod_j\cosh\frac{\pi(S-2\nu_j)}2,
 \qquad
 B_*=\frac{\pi}{4\sqrt{4\pi}\cD}
       \sqrt{\prod_j\frac{\sinh\pi\nu_j}{\nu_j}}.
\end{equation}
Expanding \eqref{eq:bd} into eight terms and using $I_0$ gives
\begin{center}
\begin{tabular}{ll}
\toprule
Conjugated legs & $T_{\eta,000}$\\
\midrule
Zero or three & $-B_*$\\
One, with conjugated leg $r$ & $B_*\cosh(\pi\nu_r)$\\
Two, with unconjugated leg $r$ & $B_*\cosh(\pi\nu_r)$\\
\bottomrule
\end{tabular}
\end{center}
For annihilation $r$ is the outgoing observer; for emission it is
the incoming observer. Reversing the observer legs relates the two
by Hermitian conjugation. These hyperbolic factors originate in the
BD frequency mixture.

Define the unit-modulus coefficients
\begin{equation}
 c_l(\nu)=i^l\prod_{j=1}^l\frac{-i\nu-j}{\sqrt{\nu^2+j^2}}
 =(-i)^l\frac{(1+i\nu)_l}{d_l(\nu)}.
\end{equation}
Normalized future data for a finite mode combination are
$A_\nu=\sum a_{lm}c_l(\nu)Y_{lm}$. This map preserves the KG
coefficient norm as an $L^2(S^2)$ norm. The late-time expansion is
\begin{equation}
 u(t,n)=2N_\nu e^{-t}
 \left[e^{-i\nu t}A_\nu(n)
       -q_\nu e^{i\nu t}\cS_\nu A_\nu(n)\right]+O(e^{-3t}),
\end{equation}
where
\begin{equation}
 \cS_\nu Y_{lm}=\frac{(1-i\nu)_l}{(1+i\nu)_l}Y_{lm},
 \qquad \cS_{-\nu}\cS_\nu=1.
\end{equation}
The normalized spherical shadow kernel is
$i\nu(\pi4^{i\nu})^{-1}|n-n'|^{-2+2i\nu}$, defined by distributional
analytic continuation. Its ordinary integral is not absolutely
convergent at these parameters. The two future branches are related
data in one BD space.

Put $w_j=1+i\eta_j\nu_j$ and
\begin{equation}
 a_i=-\frac12+\frac i2(\eta_j\nu_j+\eta_k\nu_k-\eta_i\nu_i),\qquad
 \mathcal K_\eta=\prod_{\rm cyc}|n_j-n_k|^{2a_i}.
 \label{eq:kernel}
\end{equation}
The trilinear kernel has weights $2-w_j$, not $w_j$.
Under stereographic projection with $\Omega(z)=2/(1+|z|^2)$,
$\dd\Omega_2=\Omega^2\dd^2z$,
$|n_i-n_j|^2=\Omega_i\Omega_j|z_i-z_j|^2$, and
$A_{\rm plane}=\Omega^w A_{\rm sphere}$.
All factors cancel since $2-w_i+a_j+a_k=0$.
For annihilation the resulting plane kernel is exactly DS (17)
\cite{ds}, with inputs $(i,p)$, output $o$, and $a_{\rm DS}=-i\nu$.
Emission uses the conjugate reverse kernel.

\newpage
\section{The normalized spherical period fixes the coefficient}
Let $A=a_1+a_2+a_3$ and normalize the spherical average by
\begin{equation}
 J(a_1,a_2,a_3)=\frac{1}{(4\pi)^3}\int(\dd\Omega)^3
 \prod_{\rm cyc}|n_j-n_k|^{2a_i}
 =4^A\frac{\Gamma(A+2)\prod_i\Gamma(a_i+1)}
             {\prod_i\Gamma(A-a_i+2)}.
 \label{eq:period}
\end{equation}
The domain $\operatorname{Re}a_i>-1$, $\operatorname{Re}A>-2$
ensures absolute convergence; our principal labels have
$\operatorname{Re}a_i=-1/2$, $\operatorname{Re}A=-3/2$.

For an independent normalization check, the spherical convolution
with $(2-2x)^a$ has eigenvalue
\begin{equation}
 \lambda_l(a)=\frac{4\pi\,4^a}{a+1}\frac{(-a)_l}{(a+2)_l}.
\end{equation}
Taking the trace of three such convolutions reduces $J$ to
\begin{equation}
 \frac{4^A}{\prod_i(a_i+1)}
 \sum_{l=0}^{\infty}(2l+1)
       \prod_i\frac{(-a_i)_l}{(a_i+2)_l}.
\end{equation}
The Rogers--Dougall sum, DLMF 16.4.9 \cite{dlmf}, yields
\eqref{eq:period}. Establish the trace identity for smooth kernels
first and continue analytically to the convergence domain.
The summand is $O(l^{-2\operatorname{Re}A-5})$, hence $O(l^{-2})$
here. Direct checks give $J(0,0,0)=1$ and
$J(a,0,0)=4^a/(a+1)$. This rederives a standard triangular integral
\cite{begmp}, not a new summation identity.

Define
$\mathcal B_\eta=\int\prod_j\dd\Omega_j\,A_j^{(\eta_j)}
\mathcal K_\eta$.
For generic principal-series labels the standard conformal
three-point covariance fixes this kernel up to one overall
coefficient, as in DS (17)--(19).
The actual integral fixes that coefficient without fitting:
\begin{equation}
 \boxed{\quad T_\eta=C_\eta\mathcal B_\eta,\qquad
 C_\eta=\frac{T_{\eta,000}}{(4\pi)^{3/2}J(a_1,a_2,a_3)}.\quad}
 \label{eq:coefficient}
\end{equation}
All normalized $l=0$ boundary data are $Y_{00}$, explaining the
factor $(4\pi)^{3/2}$. Equation \eqref{eq:coefficient} retains the
complex phase. Gamma reflection simplifies its norm to
\begin{equation}
 |C_\eta|^2=\frac{F_\eta}{64\pi^3\cD},\qquad
 F_\eta=
 \begin{cases}
 1,&\text{zero or three conjugated legs},\\
 \cosh^2(\pi\nu_r),&\text{one or two conjugated legs, $r$ as above}.
 \end{cases}
 \label{eq:norm}
\end{equation}

For the nontrivial $(l,l,0)$ entries, rotational averaging produces
$P_l(n_1\cdot n_2)$; expanding $P_l(1-|n_1-n_2|^2/2)$ gives a
finite sum of \eqref{eq:period} with integer shifts of $a_3$.
The two $c_l$ phases must also be included.
Independent bulk Beta sums match these boundary entries at
$l=1,2,3$, and separate boost Ward tests check higher angular
couplings. These tests support the normalization and conjugation
conventions, not an extension to a complete physical operator algebra.

\newpage
\section{Comparison with the paired vertex and its remaining scope}
The frozen fixed-mass paired $C_0$ factor extracted from CLP
\cite{clp}, up to its common convention-dependent constant, is
\begin{equation}
 D_{\rm CLP}=\frac{\pi e^{\pi S}}{\cD}
             \prod_j\frac{\sinh\pi\nu_j}{\nu_j}.
\end{equation}
Equations \eqref{eq:coefficient}--\eqref{eq:norm} imply
\begin{equation}
 \frac{|C_\eta|^2}{D_{\rm CLP}}
 =\frac{F_\eta}{64\pi^4}
       \prod_j\frac{\nu_j e^{-\pi\nu_j}}{\sinh\pi\nu_j}.
 \label{eq:comparison}
\end{equation}
The ratio consists of explicit leg factors, with the distinguished
$\cosh^2$ factor fixed by the annihilation/emission orientation.
This provides a local $G=0$ normalization interface. It does not
repair the earlier normalization-independent mismatch between the
recoil graph $H$ (or $Z_{\rm rec}$) and the proposed global $6j$.

Nor can \eqref{eq:comparison} simply be substituted for the unknown
two-point insertion $n$. That spectrum also involved the observer
weight $z$, a mass integration protocol, state preparation and
sewing norms. Choosing \eqref{eq:insertion} fixes a coupling in
this model; symmetry alone would still allow mass-dependent
couplings in other models.

Restoring the radius gives each KG-normalized scalar mode a factor
$\ell^{-1/2}$ and the three-mode volume integral a factor
$\ell^{3/2}$. A physical observable additionally requires the
dimensions and normalization of its chosen local coupling.
No finite-$G$ four-point dynamics is inferred from this free-field
calculation.

The immediate next task is to distinguish fixed-mass KG
normalization from spacetime $L^2$ or mass-delta normalization
of an observer. The global radial spectrum has two continuum
branches; its discrete sector and any imposed mass threshold must
also be identified before projecting to a BD observer space.
Only then should the coefficient, DS reflection term and spectral
measure be assembled into a physical two-point object.

The independent free-probe, large-observer-mass benchmark remains
\begin{equation}
 g_\nu(0)g_{\nu'}(0)=
 \frac{\nu\nu'}{16\pi^2\sinh\pi\nu\,\sinh\pi\nu'}.
\end{equation}
It is a disconnected antipodal product, not a zero-energy graph
value. Defining an actual mass-space operator, its group-averaged
domain, physical four-point sewing and finite-$G$ corrections
remains necessary for the full holographic objective.

\newpage
\section{Verification and source cautions}
The saved mode/vertex diagnostics contain 249 passed checks:
108 KG, ODE and independent Legendre checks; 48 spherical-entry
and coefficient-norm checks; 27 higher-angular boundary matches;
27 boost Ward checks; eight direct-integration/refinement checks;
six spherical-period checks; 16 shadow checks; and nine precision
comparisons. A separate record contains 24 checks of
\eqref{eq:comparison}. These counts are implementation checks,
not independent physical predictions.

Mass triples were $(.7,1.1,.6)$, $(1.4,.9,1.2)$ and $(3.2,2.7,.5)$.
Calculations used 60 digits, with selected nontrivial entries
repeated at 40 digits. The maximum recorded scaled error was
$6.34\times10^{-38}$. Four direct time integrals used $T=96$,
with subdivisions of width three and two; the largest relative
difference from the Beta result was $1.55\times10^{-41}$.
The corresponding analytic tail-to-signal bound was at most
$2.46\times10^{-40}$. The largest higher-angular cancellation
ratio was $632.41$. Interior quadratures are not interval-certified.
Python 3.14.5, mpmath 1.3.0, and SymPy 1.14.0, along
with full complex values and source hashes, are saved.

Two source inconsistencies are explicitly isolated.
MM HTML (14) prints a $\cosh^2t$ angular coefficient; the metric
(4), LMW (3) and the normalized modes instead require $\sech^2t$.
We derived the equation from the metric.
Also, combining BEGMP HTML (93), (127), (128) at $d=3$ gives
$J(0,0,0)=16$, inconsistent with their normalized definition.
The same constant appears in the cached PDF's (74), (101), (102).
We use the independent derivation \eqref{eq:period}, with its
zero- and single-edge checks. Neither observation is presented
as an author-confirmed erratum; source files remain unchanged.

\begin{thebibliography}{9}
\bibitem{lmw} P. Lagogiannis, A. Maloney and Y. Wang,
\emph{Odd-dimensional de Sitter Space is Transparent},
\href{https://arxiv.org/html/1106.2846v1}{arXiv:1106.2846v1}.
\bibitem{mm} D. Marolf and I. A. Morrison,
\emph{Group Averaging for de Sitter free fields},
\href{https://arxiv.org/html/0810.5163v1}{arXiv:0810.5163v1}.
\bibitem{ds} S. E. Derkachov and V. P. Spiridonov,
\emph{On the 6j-symbols for SL(2,C) group},
\href{https://arxiv.org/html/1711.07073v2}{arXiv:1711.07073v2}.
\bibitem{kl} D. Kolchmeyer and H. Liu,
\emph{Chaos and the Emergence of the Cosmological Horizon},
\href{https://arxiv.org/html/2411.08090v1}{arXiv:2411.08090v1}.
\bibitem{begmp} J. Bros, H. Epstein, M. Gaudin, U. Moschella and V. Pasquier,
\emph{Triangular invariants, three-point functions and particle stability
on the de Sitter universe},
\href{https://arxiv.org/html/0901.4223v1}{arXiv:0901.4223v1}.
\bibitem{dlmf} NIST Digital Library of Mathematical Functions,
\href{https://dlmf.nist.gov/16.4.E9}{Eq.\ 16.4.9, Rogers--Dougall sum}.
\bibitem{clp} J. Chandra, D. Liska and B. Post,
\emph{3D de Sitter from complex Virasoro TQFT: Observers and knots},
\href{https://arxiv.org/html/2610.30247v1}{arXiv:2610.30247v1}.
\end{thebibliography}
\end{document}
